% Options for packages loaded elsewhere
\PassOptionsToPackage{unicode}{hyperref}
\PassOptionsToPackage{hyphens}{url}
\PassOptionsToPackage{dvipsnames,svgnames,x11names}{xcolor}
\documentclass[
  11pt,
  a4paper,
]{article}
\usepackage{xcolor}
\usepackage[margin=2.6cm]{geometry}
\usepackage{amsmath,amssymb}
\setcounter{secnumdepth}{5}
\usepackage{iftex}
\ifPDFTeX
  \usepackage[T1]{fontenc}
  \usepackage[utf8]{inputenc}
  \usepackage{textcomp} % provide euro and other symbols
\else % if luatex or xetex
  \usepackage{unicode-math} % this also loads fontspec
  \defaultfontfeatures{Scale=MatchLowercase}
  \defaultfontfeatures[\rmfamily]{Ligatures=TeX,Scale=1}
\fi
\usepackage{lmodern}
\ifPDFTeX\else
  % xetex/luatex font selection
  \setmainfont[]{Libertinus Serif}
  \setmonofont[Scale=0.85]{Latin Modern Mono}
  \setmathfont[]{Latin Modern Math}
\fi
% Use upquote if available, for straight quotes in verbatim environments
\IfFileExists{upquote.sty}{\usepackage{upquote}}{}
\IfFileExists{microtype.sty}{% use microtype if available
  \usepackage[]{microtype}
  \UseMicrotypeSet[protrusion]{basicmath} % disable protrusion for tt fonts
}{}
\makeatletter
\@ifundefined{KOMAClassName}{% if non-KOMA class
  \IfFileExists{parskip.sty}{%
    \usepackage{parskip}
  }{% else
    \setlength{\parindent}{0pt}
    \setlength{\parskip}{6pt plus 2pt minus 1pt}}
}{% if KOMA class
  \KOMAoptions{parskip=half}}
\makeatother
\usepackage{longtable,booktabs,array}
\newcounter{none} % for unnumbered tables
\usepackage{calc} % for calculating minipage widths
% Correct order of tables after \paragraph or \subparagraph
\usepackage{etoolbox}
\makeatletter
\patchcmd\longtable{\par}{\if@noskipsec\mbox{}\fi\par}{}{}
\makeatother
% Allow footnotes in longtable head/foot
\IfFileExists{footnotehyper.sty}{\usepackage{footnotehyper}}{\usepackage{footnote}}
\makesavenoteenv{longtable}
\setlength{\emergencystretch}{3em} % prevent overfull lines
\providecommand{\tightlist}{%
  \setlength{\itemsep}{0pt}\setlength{\parskip}{0pt}}

\providecommand{\E}{\mathbb{E}}
\providecommand{\N}{\mathbb{N}}
\providecommand{\Z}{\mathbb{Z}}
\providecommand{\R}{\mathbb{R}}
\providecommand{\eps}{\varepsilon}
\providecommand{\ceil}[1]{\left\lceil #1\right\rceil}
\providecommand{\floor}[1]{\left\lfloor #1\right\rfloor}
\AtBeginDocument{\let\setminus\smallsetminus}
\usepackage[most]{tcolorbox}
\newtcolorbox{warning}{colback=red!5,colframe=red!65!black,boxrule=0.8pt,arc=2pt,left=6pt,right=6pt,top=5pt,bottom=5pt}
\usepackage{etoolbox}
\AtBeginEnvironment{longtable}{\small}
\setlength{\emergencystretch}{3em}
\usepackage{fancyhdr}
\pagestyle{fancy}
\fancyhf{}
\fancyhead[L]{\small\itshape Candidate result imbalance\_conjecture --- GPT-6 Astra, awaiting review}
\fancyhead[R]{\small\thepage}
\renewcommand{\headrulewidth}{0.2pt}
\usepackage{bookmark}
\IfFileExists{xurl.sty}{\usepackage{xurl}}{} % add URL line breaks if available
\urlstyle{same}
\hypersetup{
  pdftitle={A candidate proof of Imbalance conjecture},
  pdfauthor={github.com/graph-theory-AI --- machine-generated candidate writeup},
  colorlinks=true,
  linkcolor={blue!50!black},
  filecolor={Maroon},
  citecolor={Blue},
  urlcolor={blue!50!black},
  pdfcreator={LaTeX via pandoc}}

\title{A candidate proof of Imbalance conjecture}
\author{\href{https://github.com/graph-theory-AI}{github.com/graph-theory-AI}
--- machine-generated candidate writeup}
\date{16 September 2026}

\begin{document}
\maketitle
\begin{abstract}
A truncated-tail inequality for edge imbalances implies every
Erdős--Gallai inequality and proves the conjecture. This candidate proof
was generated by GPT-6 Astra in a single model pass for the Graph Theory
AI project. It was selected from the campaign because the model marked
it \texttt{would\_publish:\ true}; it has not been checked by the
adversarial referee, a human mathematician, or a novelty search.
\end{abstract}

\section{The problem and the claimed
result}\label{the-problem-and-the-claimed-result}

This document concerns
\href{https://graph-theory-ai.github.io/graph-conjectures/op/imbalance_conjecture/}{Imbalance
conjecture} from
\emph{http://www.openproblemgarden.org/op/imbalance\_conjecture}.

\subsection{Problem statement}\label{problem-statement}

Conjecture Suppose that for all edges \(e\in E(G)\) we have \(imb(e)>0\)
. Then \(M_{G}\) is graphic.

\subsection{Catalog context}\label{catalog-context}

Consider simple undirected graph \(G\) and let \(e=uv\in E(G)\) . The
imbalance of the edge \(e\) defined as \(imb(e)=|deg(u)-deg(v)|\) . The
multiset of all edge imbalances of \(G\) is denoted by \(M_{G}\) . Note,
that conjecture is verified for all such graphs with \(\leq 9\)
vertices.

\subsection{Claimed result}\label{claimed-result}

\textbf{Model verdict: proved} (confidence: high).

A truncated-tail inequality for edge imbalances implies every
Erdős--Gallai inequality and proves the conjecture.

\emph{The model's own caveat:} The argument is self-contained apart from
the standard Erdős--Gallai criterion; publication novelty has not been
checked.

\section{Model writeup}\label{model-writeup}

The text below is the model output verbatim, apart from moving its
machine-readable verdict block into the summary above and shifting
Markdown heading levels for this document. Mechanical TeX defects and
equation tags were normalized where needed for compilation.

\subsection{Theorem}\label{theorem}

Let \(G\) be a finite simple undirected graph such that adjacent
vertices have different degrees. Then \[
M_G=\bigl\{|d(u)-d(v)|:uv\in E(G)\bigr\}
\] is the degree multiset of a simple graph.

Write \[
a_e=|d(u)-d(v)|\qquad(e=uv),\qquad
m=|E(G)|,\qquad D=\Delta(G).
\] The hypothesis is \(a_e\ge 1\) for every edge.

The main ingredient is the following inequality.

\subsection{Truncated-tail lemma}\label{truncated-tail-lemma}

If \(m>0\), then for every integer \(1\le k<D\) and every set
\(S\subseteq E(G)\) of \(k\) edges, \[
\boxed{\quad
\sum_{e\in E(G)\setminus S}\min\{k,a_e\}\ \ge\ k(D-k).
\quad} \qquad\text{(1)}
\]

\subsubsection{Proof}\label{proof}

First, \(m\ge D\), and the hypothesis implies \(D\ge2\).

For \(k=1\), positivity gives \[
\sum_{e\notin S}\min\{1,a_e\}=m-1\ge D-1,
\] as required.

We also dispose of \(D=3,\ k=2\). If \(m\ge4\), positivity gives a lower
bound \(m-2\ge2\). If \(m=3\), the degree-three vertex is incident with
every edge, so \(G\) consists of \(K_{1,3}\) and isolated vertices. All
three imbalances are \(2\), and (1) again holds. These observations
cover all cases with \(D\le3\).

Hence assume \(D\ge4\) and \(2\le k<D\).

Fix a vertex \(v\) of degree \(D\). Let \[
q=m-D
\] be the number of edges not incident with \(v\), and put \[
s=|S\cap E(v)|,\qquad t=|S\setminus E(v)|=k-s.
\] Define \[
U=\{u\in N(v):vu\notin S\},\qquad r=|U|=D-s.
\] Finally, put \[
b=D-k.
\] Thus \(b\ge1\), \(0\le t\le k\), and \[
r=b+t. \qquad\text{(2)}
\]

There are exactly \(q-t\) edges outside \(S\) that are not incident with
\(v\). Each contributes at least \(1\) to the left side of (1).
Moreover, since \(v\) has maximum degree, \[
a_{vu}=D-d(u).
\] Consequently, \begin{equation*}
\begin{aligned}
\sum_{e\notin S}\min\{k,a_e\}
&\ge q-t+\sum_{u\in U}\min\{k,D-d(u)\}\\
&=q-t+kr-\sum_{u\in U}(d(u)-b)_+,
\end{aligned} \qquad\text{(3)}
\end{equation*} where \(x_+=\max\{x,0\}\).

Let \[
L=\{u\in U:d(u)>b\},\qquad \ell=|L|.
\] Because every vertex of \(L\) is adjacent to \(v\), \[
\sum_{u\in L}(d(u)-1)
\] counts incidences of vertices in \(L\) with edges not incident with
\(v\). Such an edge is counted at most once, except that an edge with
both endpoints in \(L\) is counted twice. Simplicity therefore gives \[
\sum_{u\in L}(d(u)-1)
\le q+|E(G[L])|
\le q+\binom{\ell}{2}.
\] It follows that \[
\sum_{u\in U}(d(u)-b)_+
=\sum_{u\in L}(d(u)-b)
\le q+\binom{\ell}{2}-(b-1)\ell. \qquad\text{(4)}
\]

Set \[
F_b(\ell)=\binom{\ell}{2}-(b-1)\ell.
\] Combining (3) and (4), we obtain \[
\sum_{e\notin S}\min\{k,a_e\}
\ge kr-t-F_b(\ell). \qquad\text{(5)}
\]

We claim that \[
F_b(\ell)\le (k-1)t. \qquad\text{(6)}
\] Indeed, \(0\le\ell\le r=b+t\), and \(F_b(x)\) is a convex quadratic.
Its maximum on this interval is attained at an endpoint. Hence
\begin{equation*}
\begin{aligned}
F_b(\ell)
&\le \max\{F_b(0),F_b(b+t)\}\\
&=\max\left\{0,\binom{t+1}{2}-\binom b2\right\}. \qquad\text{(7)}
\end{aligned}
\end{equation*}

If \(k\ge3\), then \[
\max\left\{0,\binom{t+1}{2}-\binom b2\right\}
\le \binom{t+1}{2}
\le \frac{t(k+1)}2
\le (k-1)t.
\] If \(k=2\), then \(b=D-2\ge2\) and \(t\in\{0,1,2\}\). The last
expression in (7) is respectively at most \(0,0,2\), and hence is at
most \(t=(k-1)t\). This proves (6).

Using (6) in (5), and then (2), gives \[
\sum_{e\notin S}\min\{k,a_e\}
\ge kr-t-(k-1)t
=k(r-t)
=kb
=k(D-k).
\] This proves the lemma. \(\square\)

\subsection{Completion of the proof}\label{completion-of-the-proof}

If \(m=0\), the multiset is empty and is graphic. Assume \(m>0\).

First, the imbalance sum is even: \begin{equation*}
\begin{aligned}
\sum_{uv\in E(G)}|d(u)-d(v)|
&\equiv \sum_{uv\in E(G)}(d(u)+d(v))\\
&=\sum_{v\in V(G)}d(v)^2\\
&\equiv\sum_{v\in V(G)}d(v)
=2m
\equiv0\pmod2.
\end{aligned} \qquad\text{(8)}
\end{equation*}

Also, both endpoints of an edge have positive degree, so \[
1\le a_e\le D-1. \qquad\text{(9)}
\]

Order the imbalances as \[
a_1\ge a_2\ge\cdots\ge a_m.
\] By the Erdős--Gallai criterion, it remains to prove, for every
\(1\le k\le m\), \[
\sum_{i=1}^{k}a_i
\le k(k-1)+\sum_{i=k+1}^{m}\min\{k,a_i\}. \qquad\text{(10)}
\]

If \(k\ge D\), then (9) gives \[
\sum_{i=1}^{k}a_i
\le k(D-1)
\le k(k-1),
\] so (10) holds.

If \(1\le k<D\), apply the truncated-tail lemma to a set of edges
carrying the \(k\) largest imbalances. It gives \[
\sum_{i=k+1}^{m}\min\{k,a_i\}\ge k(D-k).
\] Therefore, \begin{equation*}
\begin{aligned}
k(k-1)+\sum_{i=k+1}^{m}\min\{k,a_i\}
&\ge k(k-1)+k(D-k)\\
&=k(D-1)\\
&\ge\sum_{i=1}^{k}a_i.
\end{aligned}
\end{equation*} Thus every Erdős--Gallai inequality holds. Together with
(8), this proves that \(M_G\) is graphic. \(\square\)

\subsection{Scope and verification
status}\label{scope-and-verification-status}

The argument covers disconnected graphs, isolated vertices, the empty
graph, and all maximum degrees. Simplicity is used explicitly in the
estimate \[
|E(G[L])|\le\binom{|L|}{2},
\] and positive imbalance is used to obtain one unit of tail
contribution from every retained edge not incident with the chosen
maximum-degree vertex.

No conjectural auxiliary assertion or computational check is needed. I
have not independently checked whether this proof has appeared in the
literature.

\appendix

\section{Provenance and review
status}\label{provenance-and-review-status}

\begin{warning}

\textbf{UNREFEREED MODEL OUTPUT.} This document was selected solely
because GPT-6 Astra labelled its own result
\texttt{would\_publish:\ true} and returned \texttt{proved} or
\texttt{disproved}. It has not passed the adversarial LLM referee used
for the earlier Sol campaign, has not been checked by a human
mathematician, and has not been checked for novelty. Treat every
mathematical and bibliographic claim as unverified.

\end{warning}

{\def\LTcaptype{none} % do not increment counter
\begin{longtable}[]{@{}
  >{\raggedright\arraybackslash}p{(\linewidth - 2\tabcolsep) * \real{0.5000}}
  >{\raggedright\arraybackslash}p{(\linewidth - 2\tabcolsep) * \real{0.5000}}@{}}
\toprule\noalign{}
\endhead
\bottomrule\noalign{}
\endlastfoot
Catalog id & \texttt{imbalance\_conjecture} \\
Catalog entry &
\href{https://graph-theory-ai.github.io/graph-conjectures/op/imbalance_conjecture/}{Imbalance
conjecture} \\
Source corpus & opg \\
Campaign leg & OpenProblemGarden first pass (\texttt{attacks\_opg}) \\
Source paper / entry &
http://www.openproblemgarden.org/op/imbalance\_conjecture \\
Model verdict & \textbf{proved} (confidence: high) \\
Selection criterion & \texttt{would\_publish:\ true} and verdict
\texttt{proved} or \texttt{disproved} \\
Model & \texttt{gpt-6-astra}, reasoning effort \texttt{max},
\texttt{mode=pro}, flex service tier \\
Original artifact &
\texttt{attacks\_opg/imbalance\_conjecture/output.md} \\
Independent review & \textbf{None} \\
Model's caveats & The argument is self-contained apart from the standard
Erdős--Gallai criterion; publication novelty has not been checked. \\
\end{longtable}
}

The generated Markdown and LaTeX sources for this PDF are stored under
\texttt{to\_review\_astra/src/imbalance\_conjecture/}.

\end{document}
